\documentclass[11pt]{article}
\usepackage[margin=2.5cm]{geometry}
\usepackage{amsmath,amssymb,amsthm}
\usepackage{enumitem}
\usepackage{verbatim}
\usepackage{url}

\newtheorem{theorem}{Theorem}
\newtheorem{lemma}[theorem]{Lemma}
\newtheorem{definition}[theorem]{Definition}
\newtheorem{remark}[theorem]{Remark}

\newcommand{\Deriv}{\mathsf{Deriv}}
\newcommand{\thmT}{\mathsf{thmT}}
\newcommand{\Con}{\mathsf{Con}}
\newcommand{\bot}{\mathsf{bot}}
\newcommand{\Tree}{\mathsf{Tree}}
\newcommand{\Term}{\mathsf{Term}}
\newcommand{\FunOne}{\mathsf{Fun_1}}
\newcommand{\FunTwo}{\mathsf{Fun_2}}
\newcommand{\Formula}{\mathsf{Formula}}
\newcommand{\Pair}{\mathsf{Pair}}
\newcommand{\Fst}{\mathsf{Fst}}
\newcommand{\Snd}{\mathsf{Snd}}
\newcommand{\IfLf}{\mathsf{IfLf}}
\newcommand{\Lift}{\mathsf{Lift}}
\newcommand{\Comp}{\mathsf{Comp}}
\newcommand{\Fan}{\mathsf{Fan}}
\newcommand{\Post}{\mathsf{Post}}
\newcommand{\KT}{\mathsf{KT}}
\newcommand{\reify}{\mathsf{reify}}
\newcommand{\code}{\mathsf{code}}
\newcommand{\codeFormula}{\mathsf{codeFormula}}
\newcommand{\codeTriv}{\mathsf{codeTriv}}
\newcommand{\falseT}{\mathsf{falseT}}
\newcommand{\subT}{\mathsf{subT}}
\newcommand{\subTT}{\mathsf{subT_2}}
\newcommand{\imp}{\Rightarrow}

\title{G\"odel's Second Incompleteness Theorem\\
  for Binary Recursive Arithmetic in Agda}
\author{Thierry Coquand}
\date{May 2026}

\begin{document}
\maketitle

\begin{abstract}
We formalize G\"odel's second incompleteness theorem (G\"odel~II) for
\emph{Binary Recursive Arithmetic} (BRA), a Rose/Guard-style equational
system on binary trees with explicit substitution and primitive
recursion.  The development is in Agda with
\verb|--safe --without-K --exact-split|, ASCII-only, no postulates, no
holes.  The headline result is
\(\thmT \text{ is sound, and } \Deriv\,\Con \to \Deriv\,\bot.\)
The proof is constructive and consists of an internalised proof
checker $\thmT$ (Theorem~11), a Theorem~12 closure expressing that all
primitive-recursive functions of BRA are reflected in $\thmT$, a
Theorem~14 closure linking $\thmT$ to BRA's encoding of provability,
and a final Compose module producing $\Deriv\,\Con \to \Deriv\,\bot$.
\end{abstract}

%% ============================================================
\section{Overview}
%% ============================================================

The BRA codebase lives in \texttt{BRA/}; this document summarises its
main result.  The chain is

\[
\text{Thm 11 (concrete)} \;\to\;
\text{Thm 12 (parametric)} \;\to\;
\text{Thm 13} \;\to\;
\text{Thm 14} \;\to\;
\text{Compose} \;\to\;
\text{godelII}
\]

assembled in
\verb|BRA/GoedelIIFull.agda|, whose top-level statement is
\begin{verbatim}
godelII : Deriv Con -> Deriv bot
\end{verbatim}

After the soundification work described in
Section~\ref{sec:sound}, the meaning of $\Con$ is robust: every
output of $\thmT$ on a syntactically well-formed proof tree is either
a valid encoded formula $\codeFormula(P)$ or the explicit safe default
$\codeTriv = \falseT = \codeFormula(0=0)$, so $\Con \equiv \neg
\Pr(\codeFormula(\bot))$ truly expresses BRA's consistency.

%% ============================================================
\section{The base system}
%% ============================================================

\subsection{Trees, terms, formulas}

\begin{definition}[Tree]
\(\Tree ::= \mathsf{lf} \mid \mathsf{nd}\;\Tree\;\Tree.\)
\end{definition}

\begin{definition}[Term, Functors]
$\Term$ has constructors $\mathsf{O}$, $\mathsf{var}\;n$ for $n :
\mathbb{N}$, $\mathsf{ap_1}\;f\;t$ for $f : \FunOne$, and
$\mathsf{ap_2}\;g\;t_1\;t_2$ for $g : \FunTwo$.  $\FunOne$ and
$\FunTwo$ are inductive types: $\Fst, \Snd, \Lift\,g, \Comp\,f\,g,
\KT\,t \in \FunOne$; $\Pair, \IfLf, \Fan\,f\,g\,h, \Post\,f\,g \in
\FunTwo$, etc.
\end{definition}

\begin{definition}[Formula]
$\Formula ::= \mathsf{atomic}\,(\mathsf{eqn}\,t_l\,t_r) \mid
\mathsf{not}\,P \mid P \mathbin{\mathsf{imp}} Q.$
\end{definition}

\subsection{Derivability}

$\Deriv : \Formula \to \mathsf{Set}$ is an inductive predicate with
explicit constructors for the BRA axioms (axI, axFst, axSnd,
axComp, axComp2, axLift, axPost, axFan, axKT, axIfLfL, axIfLfN,
axTreeEq[LL/LN/NL/NN], axGoodstein, axRefl, axK, axS, axNeg,
axContrapos, axExFalso, axEqTrans, axEqCong[1/L/R], axRecLf, axRecNd,
axRecPLf, axRecPNd, axFstLf, axSndLf, axIfLfLL, axIfLfNL, axConst)
and the inductive rules
$\mathsf{ruleSym},\,\mathsf{ruleTrans},\,\mathsf{cong_1},\,\mathsf{cong_L},\,
\mathsf{cong_R},\,\mathsf{mp},\,\mathsf{ruleInst},\,\mathsf{ruleInst_2},\,
\mathsf{ruleIndBT}$.

\subsection{Encoding and reify}

Each $\Term$ has a tree code $\code : \Term \to \Tree$; each
$\Formula$ has a code $\codeFormula : \Formula \to \Tree$.  Trees lift
to terms via $\reify : \Tree \to \Term$.  The compositions
$\reify \circ \code$ and $\code \circ \reify$ are not the identity in
general, but $\reify \circ \codeFormula\,P$ is always Pair-shaped
because $\codeFormula(P) = \mathsf{nd}\;\_\;\_$ for every $P$.

%% ============================================================
\section{Theorem 11: the proof checker $\thmT$}
%% ============================================================

\begin{definition}[\thmT]
$\thmT : \FunOne$ is a one-argument BRA function, defined inside the
$\mathsf{abstract}$ block of \verb|BRA/Thm/ThmT.agda| via a
treeRec--style fold over its proof-tree input.  At each tag
$\mathsf{tagCode}_X$, $\thmT$ dispatches to a body $\mathsf{body}_X :
\FunTwo$ that computes the equation code from the proof's payload and
sub-proof images.
\end{definition}

\begin{theorem}[Thm 11, the dispatchers]
For every BRA derivation rule $X$ with encoder $\mathsf{enc}_X$ and
result $\mathsf{out}_X$, there is a closed-form lemma
\verb|thmTDispX| in \verb|BRA/Thm/ThmT.agda| of the shape
\[
  \forall\; (\text{IH-equations}),\;
  \Deriv\big((\mathsf{ap_1}\;\thmT\;\reify(\mathsf{enc}_X\,\ldots)) =
            \reify(\mathsf{out}_X\,\ldots)\big).
\]
Each $\mathsf{thmTDispX}$ proves that $\thmT$ at the encoded proof
tree reduces to the encoded conclusion, given the IHs supplied for
the sub-proof trees.
\end{theorem}

%% ============================================================
\section{Soundification of $\thmT$}
\label{sec:sound}
%% ============================================================

The original $\mathsf{body}_X$ definitions (modelled on Guard
section~3.3) are \emph{unconditional}: they extract sub-proof images
from the input pair $(a,bb)$ via $\mathsf{Lift}/\mathsf{Comp}/\mathsf{Post}$
combinators and assemble the conclusion.  For a malformed input proof
tree, the unconditional bodies can produce arbitrary terms via
safe-default fallbacks (\verb|axFstLf|, \verb|axSndLf|), risking
$\thmT(y_{\text{bad}}) = \codeFormula(\bot)$ for some adversarial
$y_{\text{bad}}$, which would refute $\Con$ in the meta-theory and
make $\mathsf{godelII}$ vacuous.

\subsection{Verifying-body design}

For each premise-consuming dispatcher $X$, we replace
$\mathsf{body}_X$ with a verifying variant
$\mathsf{body}_X^v$ of the form
\[
\mathsf{body}_X^v \;=\; \Post\;\mathsf{verifier}_X\;\Pair,
\]
where $\mathsf{verifier}_X$ is a $\FunOne$ over $t = \Pair\,a\,bb$:
\[
\mathsf{verifier}_X
\;=\;\Comp_2\;\IfLf\;\langle\text{discriminant}\rangle\;
              (\Comp_2\;\Pair\;(\KT\;\codeTriv)\;\langle\text{assembly}\rangle).
\]
The discriminant is whichever sub-proof image (or encoding-level
projection) the body would otherwise unsafely use.  On the leaf
($\mathsf{O}$) branch the body outputs $\codeTriv = \falseT$; on the
Pair branch it executes the original assembly.  The eval-pass lemma
$\mathsf{body}_X^v\!.\mathsf{eval\_pass}$ takes a Pair-shape
witness for the IH and discharges the IfLf via $\mathsf{axIfLfN}$.

\subsection{Coverage}

\begin{center}
\renewcommand{\arraystretch}{1.05}
\begin{tabular}{lll}
\hline
Dispatcher & Proto / VProof file & Variants delivered \\
\hline
mp           & \verb|BRA/SoundMpProto.agda|         & closed + lifted \\
ruleInst     & \verb|BRA/SoundRuleInstProto.agda|   & closed + lifted \\
ruleSym      & \verb|BRA/SoundRuleSymProto.agda|    & closed + lifted \\
cong$_1$     & \verb|BRA/SoundCong1Proto.agda|      & closed \\
cong$_L$     & \verb|BRA/SoundCongLProto.agda|      & closed + lifted + dl \\
cong$_R$     & \verb|BRA/SoundCongRProto.agda|     & closed + lifted + dl \\
ruleTrans    & \verb|BRA/SoundRuleTransProto.agda|  & closed + lifted + dl \\
ruleInst$_2$ & \verb|BRA/SoundRuleInst2Proto.agda|  & closed + parametric \\
ruleIndBT    & \verb|BRA/SoundRuleIndBTProto.agda|  & closed + closed-param \\
\hline
\end{tabular}
\end{center}

For each row the body in \verb|BRA/Thm/ThmT.agda| reads
$\mathsf{body}_X = \mathsf{body}_X^v$ and the dispatchers
$\mathsf{thmTDisp}X$, $\mathsf{thmTDisp}X\_\mathsf{param}$,
\ldots\ route through the corresponding $\mathsf{body}_X^v\!.\mathsf{eval\_pass}$
variants supplied with the appropriate Pair-shape witness.

\begin{theorem}[Soundness of $\thmT$]
\label{thm:sound}
For every closed proof tree $y$, $\thmT(y)$ is either $\reify
\circ \codeFormula(P)$ for some $P$, or the explicit safe default
$\codeTriv = \codeFormula(0=0)$.  In particular, $\thmT(y)
\neq \codeFormula(\bot)$ for any $y$ unless BRA is actually
inconsistent.
\end{theorem}

\begin{remark}[Lifted variants]
Several Theorem~12 schematic proofs (Parts/Comp2, Parts/Fan, etc.)
require Pair-shape hypotheses inside $P \imp \cdot$ contexts (singly or
doubly lifted).  We provide
$\mathsf{body}_X^v\!.\mathsf{eval\_pass\_l}$ and
$\mathsf{body}_X^v\!.\mathsf{eval\_pass\_dl}$ for cong$_L$, cong$_R$,
ruleTrans, mp, ruleInst, ruleSym; the closed pass suffices for cong$_1$
and ruleIndBT.
\end{remark}

%% ============================================================
\section{Theorem 12: parametric closure}
%% ============================================================

For each primitive-recursive function $f$ definable in BRA, Theorem~12
internalises the equation $\thmT(\mathsf{D}_f(x)) = \mathsf{th}(f(x))$
inside BRA itself.  The implementation uses 15 modular pairs of
``Param'' wrappers (\verb|BRA/Thm12/Param/AxX.agda|) and ``Parts''
glue (\verb|BRA/Thm12/Parts/X.agda|), assembled into the
\verb|FromBridges| record in \verb|BRA/Thm12.agda|.

\begin{theorem}[Thm 12 result records]
The records $\mathsf{Thm12\_F_1\_Result}\;\thmT$ and
$\mathsf{Thm12\_F_2\_Result}\;\mathsf{sub}$ in
\verb|BRA/Thm12.agda| are inhabited unconditionally by
$\mathsf{FromBridges.thm12\_F_1}$ and
$\mathsf{FromBridges.thm12\_F_2}$.
\end{theorem}

%% ============================================================
\section{Theorem 14 and the Compose module}
%% ============================================================

Theorem~14 packages the steps 1--5 of Guard's section~3.5 into a
contract on a single function $\mathsf{constr_5} : \FunOne$ together
with a parametric internal implication
\[
\mathsf{step_5}(x) :
  \Deriv\Big((\thmT(x) = c_G) \imp (\thmT(\mathsf{constr_5}\,x) = \codeFormula(\bot))\Big)
\]
where $c_G = \reify(\codeFormula(G))$ is the code of the G\"odel
sentence.  The concrete $\mathsf{constr_5\_final}$ and
$\mathsf{step_5\_l}$ live in \verb|BRA/Th14Step5.agda|.  The Compose
module (\verb|BRA/GoedelII.agda|) then produces

\begin{theorem}[godelII]
\label{thm:godelII}
$\Deriv\,\Con \to \Deriv\,\bot$,
where $\Con = \neg\,(\mathsf{ap_1}\,\thmT\,(\mathsf{var}\,0) =
\reify(\codeFormula(\bot)))$ universally closed in $\mathsf{var}\,0$.
\end{theorem}

The headline file \verb|BRA/GoedelIIFull.agda| instantiates Compose with
$\mathsf{constr_5\_final}$ and $\mathsf{step_5\_l}$, plugs in
$\mathsf{r_{12}\_th} = \mathsf{FromBridges.thm12\_F_1}\,\thmT$ and
$\mathsf{r_{12}\_sub} = \mathsf{FromBridges.thm12\_F_2}\,\mathsf{sub}$,
and exports

\begin{verbatim}
godelII : Deriv Con -> Deriv bot
godelII = Final.godelII
\end{verbatim}

unconditionally.

%% ============================================================
\section{Implementation notes}
%% ============================================================

\begin{itemize}
\item \textbf{Build}.  Cold rebuild of \verb|BRA/GoedelIIFull.agda|
  takes \(\sim\!32\)s.  Cached typechecks of the headline file and
  \verb|BRA/Thm/ThmT.agda| are under 1.2\,s each.

\item \textbf{Discipline}.  Throughout BRA we use only
  \verb|--safe --without-K --exact-split|.  No postulates, no holes,
  no \verb|with|-abstraction, no dot patterns.  Every let-binding uses
  camelCase to avoid Agda's mid-identifier underscore mixfix grammar.

\item \textbf{Soundification scope}.  All 9 inductive
  premise-consuming dispatchers in \verb|BRA/Thm/ThmT.agda| now have
  $\mathsf{body}_X = \mathsf{body}_X^v$.  No further dispatcher reads
  sub-proof images without an IfLf gate.

\item \textbf{Encoding-level safety}.  \verb|ruleIndBT|'s body extracts
  $\codeFormula(P)$ directly from the encoding (no IH consumed); its
  verifying variant gates on the encoding's outer $\Snd$ being
  Pair-shaped, so a forged leaf payload is mapped to $\codeTriv$.
\end{itemize}

%% ============================================================
\section{Files}
%% ============================================================

The complete BRA codebase is in the \verb|BRA/| subdirectory of the
project repository.  Headline files:

\begin{itemize}
\item \verb|BRA/GoedelIIFull.agda| --- top-level
  $\mathsf{godelII} : \Deriv\,\Con \to \Deriv\,\bot$.
\item \verb|BRA/GoedelII.agda| --- Compose module taking the
  Theorem~14 contract and producing $\mathsf{godelII}$.
\item \verb|BRA/Thm14Abstract.agda| --- abstract Theorem~14 tower.
\item \verb|BRA/Thm14Step5.agda| --- concrete
  $\mathsf{constr_5\_final}$ + $\mathsf{step_5\_l}$.
\item \verb|BRA/Thm12.agda|, \verb|BRA/Thm12/| --- Theorem~12 closure
  (15 Param + Parts pairs).
\item \verb|BRA/Thm/ThmT.agda| --- $\thmT$ definition and all
  $\mathsf{thmTDisp}X$ dispatchers (the abstract block).
\item \verb|BRA/Sound*Proto.agda|,
      \verb|BRA/Sound*VProof.agda| --- verifying bodies and their
  eval-pass lemmas (Section~\ref{sec:sound}).
\end{itemize}

\end{document}
